\section{Introduction}
\label{sec:intro}

Diet-related chronic disease accounts for a large share of preventable mortality worldwide, and smartphone cameras have made it easy to photograph meals in situ.
The harder part is converting that photograph into something actionable: a quantified answer to whether the plate in front of you resembles what nutritional guidelines actually recommend.
A growing body of computer-vision work addresses pieces of this problem---food classification~\cite{bossard2014food101}, detection~\cite{matsuda2012uecfood}, and segmentation~\cite{wu2021foodseg103}---but most published pipelines stop at identification and leave the gap between ``broccoli, rice, chicken'' and ``was this a good meal?'' entirely to the user.

\textbf{SmartPlate.}
We close that gap with a lightweight end-to-end system.
Given a single plate photo, we (i) predict a pixel-level segmentation across $103$ fine-grained food categories using a SegFormer-B0~\cite{xie2021segformer} fine-tuned on FoodSeg103~\cite{wu2021foodseg103}, (ii) aggregate the per-pixel predictions into $7$ food groups from the Harvard Healthy Eating Plate~\cite{harvardplate}, and (iii) score the plate on a $0$--$100$ scale via a transparent, rule-based formula that rewards meeting the Harvard target proportions and penalizes refined grains, red or processed meat, and sugary or fatty items.
The same pixel-area proportions drive a coarse macro breakdown through a per-group USDA~\cite{usdafdc} lookup table.

A key design choice is keeping the scoring module closed-form rather than learned: the score is a short function of food-group proportions, so the reason a plate receives a given score is immediately readable from the formula.
This confines the learned component to the part of the problem where it adds clear value---pixel-level food localization---and delegates dietary judgment to a published rubric that can be inspected and critiqued on its own terms.

\textbf{Contributions.}
\begin{enumerate}
    \item An end-to-end pipeline from a raw plate image to a Harvard-aligned health score and coarse macro breakdown, runnable on a single Apple Silicon laptop.
    \item A fine-tuned SegFormer-B0 food-segmentation model with quantitative evaluation at both $103$-class and $7$-group granularity on the FoodSeg103 validation split.
    \item A ResNet-50~\cite{he2016resnet} multi-label baseline trained on the same data, giving a controlled comparison of what pixel-level outputs buy over image-level outputs for this task.
    \item An interpretable, closed-form scoring function and a hand-built mapping from the $103$ FoodSeg103 classes to the $7$ Harvard food groups, released alongside the code.
\end{enumerate}
